\documentclass[11pt,fontset=fandol]{ctexart}
\usepackage[a4paper,margin=23mm]{geometry}
\usepackage{amsmath,amssymb,mathtools,bm,mathrsfs,esint}
\usepackage{graphicx,longtable,booktabs,array,calc}
\usepackage[normalem]{ulem}
\usepackage{xcolor,hyperref}
\hypersetup{hidelinks}
\setcounter{secnumdepth}{-1}
\newcounter{none}
\providecommand{\tightlist}{\setlength{\itemsep}{0pt}\setlength{\parskip}{0pt}}
\providecommand{\passthrough}[1]{#1}
\usepackage{amssymb,mathrsfs,esint,bm,newunicodechar}
\allowdisplaybreaks
\setlength{\emergencystretch}{3em}
\providecommand{\llangle}{\langle\!\langle}
\providecommand{\rrangle}{\rangle\!\rangle}
\newunicodechar{∞}{\ensuremath{\infty}}
\newunicodechar{∈}{\ensuremath{\in}}
\newunicodechar{∑}{\ensuremath{\sum}}
\newunicodechar{∫}{\ensuremath{\int}}
\newunicodechar{≤}{\ensuremath{\leq}}
\newunicodechar{≥}{\ensuremath{\geq}}
\newunicodechar{≠}{\ensuremath{\neq}}
\newunicodechar{→}{\ensuremath{\to}}
\newunicodechar{⇒}{\ensuremath{\Rightarrow}}
\newunicodechar{α}{\ensuremath{\alpha}}
\newunicodechar{β}{\ensuremath{\beta}}
\newunicodechar{γ}{\ensuremath{\gamma}}
\newunicodechar{λ}{\ensuremath{\lambda}}
\newunicodechar{φ}{\ensuremath{\varphi}}
\newunicodechar{ε}{\ensuremath{\varepsilon}}
\newunicodechar{σ}{\ensuremath{\sigma}}
\newunicodechar{θ}{\ensuremath{\theta}}
\newunicodechar{δ}{\ensuremath{\delta}}
\newunicodechar{η}{\ensuremath{\eta}}
\newunicodechar{π}{\ensuremath{\pi}}
\newunicodechar{ξ}{\ensuremath{\xi}}
\newunicodechar{Σ}{\ensuremath{\Sigma}}
\newunicodechar{Δ}{\ensuremath{\Delta}}
\newunicodechar{∂}{\ensuremath{\partial}}
\newunicodechar{∀}{\ensuremath{\forall}}
\newunicodechar{∃}{\ensuremath{\exists}}
\newunicodechar{√}{\ensuremath{\surd}}
\newunicodechar{−}{\ensuremath{-}}
\newunicodechar{μ}{\ensuremath{\mu}}
\newunicodechar{ρ}{\ensuremath{\rho}}
\newunicodechar{τ}{\ensuremath{\tau}}
\newunicodechar{ω}{\ensuremath{\omega}}
\newunicodechar{Ω}{\ensuremath{\Omega}}
\newunicodechar{χ}{\ensuremath{\chi}}
\newunicodechar{ψ}{\ensuremath{\psi}}
\newunicodechar{Γ}{\ensuremath{\Gamma}}
\newunicodechar{Φ}{\ensuremath{\Phi}}
\newunicodechar{Λ}{\ensuremath{\Lambda}}
\newunicodechar{ℝ}{\ensuremath{\mathbb{R}}}
\newunicodechar{ℤ}{\ensuremath{\mathbb{Z}}}
\newunicodechar{ℕ}{\ensuremath{\mathbb{N}}}
\newunicodechar{△}{\ensuremath{\triangle}}
\newunicodechar{⊗}{\ensuremath{\otimes}}
\begin{document}
\section{2014年杭州电子科技大学601数学分析考研真题}\label{ux5e74ux676dux5ddeux7535ux5b50ux79d1ux6280ux5927ux5b66601ux6570ux5b66ux5206ux6790ux8003ux7814ux771fux9898}

来源：\texttt{HDU/601数学分析/2014/2014年杭州电子科技大学601数学分析考研真题.pdf}。原
PDF 页序：1, 2, 3。

\section{杭州电子科技大学}\label{ux676dux5ddeux7535ux5b50ux79d1ux6280ux5927ux5b66}

\section{2014年攻读硕士学位研究生入学考试}\label{ux5e74ux653bux8bfbux7855ux58ebux5b66ux4f4dux7814ux7a76ux751fux5165ux5b66ux8003ux8bd5}

\section{《数学分析》试题}\label{ux6570ux5b66ux5206ux6790ux8bd5ux9898}

\begin{quote}
助手整理：修正分段函数的数组列及右极限符号。末题关于``任意一点''的截距结论，在三个坐标均为正、三个截距有定义时成立，详见整理后的参考解析；此处保留原题文字。
\end{quote}

（试题共四大题，共3页，总分150分）

姓名 报考专业 准考证号

【所有答案必须写在答题纸上，做在试卷或草稿纸上无效】

\subsection{一、填空题（本大题共8小题，每小题5分，本大题共40分）}\label{ux4e00ux586bux7a7aux9898ux672cux5927ux9898ux51718ux5c0fux9898ux6bcfux5c0fux98985ux5206ux672cux5927ux9898ux517140ux5206}

1．求极限 \(\lim _ { x \to 0 } ( 1 + \tan x ) ^ { \cot x }\)

2.设
\(f ( x ) = \left\{ \begin{array} { l l } { { 7 e ^ { x } - \cos x \mathrm { ~ , ~ } x \le 0 } } \\ { { \displaystyle { \frac { \sin 2 x } { \tan a x } } , ~ x > 0 } } \end{array} \right.\)
，要使f(x)在x=0处连续，则a=

3.设f（x)在{[}0,+∞）上连续，且
\(\int _ { 0 } ^ { x } f ( t ) d t = x ( 1 + \cos x )\) ，则f(π/2)=

4.积分
\(\int _ { - 2 } ^ { 2 } \left( \left| x \right| + x \right) e ^ { \left| x \right| } d x =\)

5.已知f（x)二阶可导，且 \(y = e ^ { f ( x ) }\) ，则
\(y ^ { \prime \prime } =\)

6.设f（x,y）是连续函数，交换二次积分
\(\int _ { 0 } ^ { 4 } d x \int _ { x } ^ { 2 \sqrt { x } } f ( x , y ) d y\)
的积分次序，可表示为

7.幂级数
\(\sum _ { n = 1 } ^ { \infty } { \frac { ( x - 3 ) ^ { n } } { n \cdot 3 ^ { n } } }\)
的收敛域为

8.设函数y=y（x）由方程 \(y - x e ^ { y } = 1\) 所确定，求
\(\left. { \frac { d ^ { 2 } y } { d x ^ { 2 } } } \right| _ { x = 0 }\)

\subsection{二、计算题（本大题共4小题，本大题共30分）}\label{ux4e8cux8ba1ux7b97ux9898ux672cux5927ux9898ux51714ux5c0fux9898ux672cux5927ux9898ux517130ux5206}

1.（本小题5分）求极限
\(\lim _ { x \to 0 } { \frac { 1 } { x } } { \left( { \frac { 1 } { \sin x } } - { \frac { 1 } { \tan x } } \right) }\)

2.（本小题5分）求积分 \(\int e ^ { e ^ { x } + x } d x\)

3.（本小题10分）设幂级数为
\(\sum \limits _ { n = 1 } ^ { \infty } n 3 ^ { n - 1 } x ^ { n - 1 }\)
，求其和函数f(x)，并求 \(\int _ { 0 } ^ { 1 / 8 } f ( x ) d x\)

4.（本小题10分）计算曲线积分
\(\int _ { . } ( e ^ { x } \sin y - m y ) d x + ( e ^ { x } \cos y - m ) d y\)
，其中L是由点O(0,0)到A(0,a）沿 \(x = \sqrt { a y - y ^ { 2 } }\)
的右半圆周.

\subsection{三、综合题（本大题共3小题，每小题10分，本大题共30分）}\label{ux4e09ux7efcux5408ux9898ux672cux5927ux9898ux51713ux5c0fux9898ux6bcfux5c0fux989810ux5206ux672cux5927ux9898ux517130ux5206}

1．讨论函数
\(f ( x , y ) = \left\{ \begin{array} { c c } { { ( x ^ { 2 } + y ^ { 2 } ) { \sin } \frac { 1 } { x ^ { 2 } + y ^ { 2 } } , } } & { { x ^ { 2 } + y ^ { 2 } \ne 0 } } \\ { { 0 , } } & { { x ^ { 2 } + y ^ { 2 } = 0 } } \end{array} \right.\)

在点(0,0)处是否可微？在点(0,0)处的偏导是否连续？

2.讨论函数项级数
\(\sum \limits _ { n = 1 } ^ { \infty } \frac { x } { 1 + n ^ { 4 } x ^ { 2 } }\)
在 \(- \infty < x < + \infty\) 上的一致收敛性.

3.设f(t)在{[}0,1{]}上连续， \(f ( 0 ) = 0 , f ^ { \prime } ( 0 ) = 1\)

\[
F ( t ) = \int _ { x ^ { 2 } + y ^ { 2 } + z ^ { 2 } \leq t ^ { 2 } } f ( { \sqrt { x ^ { 2 } + y ^ { 2 } + z ^ { 2 } } } ) d x d y d z ,
\]

讨论极限 \(\lim_{t\to0^+}\dfrac{F(t)}{t^4}\)
是否存在，若存在，求出极限值。

\subsection{四、证明题（本大题共5小题，每小题10分，本大题共50分）}\label{ux56dbux8bc1ux660eux9898ux672cux5927ux9898ux51715ux5c0fux9898ux6bcfux5c0fux989810ux5206ux672cux5927ux9898ux517150ux5206}

1.证明不等式
\(\ln \left( 1 + { \frac { 1 } { x } } \right) > { \frac { 1 } { 1 + x } } , ( 0 < x < + \infty )\)

2.设f(x)在{[}a,b{]}上可导(a\textless b)，f(a)=f(b)，试证：存在
\(\xi \in ( a , b )\)
，使\(f ( a ) - f ( \xi ) \ = \xi \ f ^ { \prime } ( \xi )\)

3.证明 \(f ( x ) = \cos { \frac { 1 } { x } }\)
在(0,1)上连续，但非一致连续.

4.设函数f(x)在{[}a,b{]}上连续，且 \(f ( x ) > 0\) ，证明：在
\(( a , b )\) 内有且仅有一点 \(\xi\) ，使得
\(\int _ { a } ^ { \xi } f ( x ) d x = \int _ { \xi } ^ { b } [ f ( x ) ] ^ { - 1 } d x\)

5.证明曲面
\({ \sqrt { x } } + { \sqrt { y } } + { \sqrt { z } } = { \sqrt { a } } \left( a > 0 \right)\)
在任意\ldots·点处的切平面在三坐标轴上的截距之和为定值。
\end{document}
